%==============================================================
%  lecons/analyse/cosinus_remarquables.tex — Valeurs remarquables du cosinus
%==============================================================

\lecontitle{Les valeurs remarquables du cosinus}{Analyse réelle — Trigonométrie et nombres algébriques}

%=============================================================
%  § 1 — DÉFINITION ET RÉSULTATS FONDAMENTAUX
%=============================================================
\secnum{navyblue}{1}{Définition et résultats fondamentaux}

\begin{tcolorbox}[thmbox]
  \textbf{Objet.}\enspace\index{valeurs remarquables du cosinus}
  On cherche les valeurs exactes de $\cos\theta$ pour les angles « classiques »
  $\theta$, c'est-à-dire les multiples rationnels de $\pi$ dont le cosinus
  s'exprime par radicaux. Les cinq angles fondamentaux donnent le tableau~:
  \[
    \renewcommand{\arraystretch}{1.6}
    \begin{array}{c|ccccc}
      \theta & 0 & \dfrac{\pi}{6} & \dfrac{\pi}{4} & \dfrac{\pi}{3} & \dfrac{\pi}{2}\\[2pt]
      \hline
      \cos\theta & 1 & \dfrac{\sqrt 3}{2} & \dfrac{\sqrt 2}{2} & \dfrac{1}{2} & 0\\[6pt]
      \sin\theta & 0 & \dfrac{1}{2} & \dfrac{\sqrt 2}{2} & \dfrac{\sqrt 3}{2} & 1
    \end{array}
  \]

  \medskip
  \textbf{Théorème (trois principes de calcul).}
  Toute valeur du tableau, et plus généralement $\cos\!\big(\tfrac{p}{q}\pi\big)$,
  se calcule par l'un des trois leviers~:
  \begin{itemize}
    \item \textit{Géométrie élémentaire} — le triangle équilatéral fournit
      $\pi/3$ et $\pi/6$, le triangle rectangle isocèle fournit $\pi/4$,
      via le théorème de Pythagore et $\cos^2\theta+\sin^2\theta=1$.
    \item \textit{Formules d'addition, duplication, bissection} — elles
      propagent les valeurs connues vers $\pi/12$, $5\pi/12$, $\pi/8$, $\pi/{24}$, \dots
      \[
        \cos(a\pm b)=\cos a\cos b\mp\sin a\sin b,\qquad
        \cos\frac{\theta}{2}=\pm\sqrt{\frac{1+\cos\theta}{2}}.
      \]
    \item \textit{Racines d'un polynôme} — pour les angles que les deux leviers
      précédents n'atteignent pas (typiquement $\pi/5$), on écrit une relation
      polynomiale issue de $\cos(n\theta)$ et on la résout lorsque c'est possible.
      Ainsi
      \[
        \cos\frac{\pi}{5}=\frac{1+\sqrt 5}{4},\qquad
        \cos\frac{2\pi}{5}=\frac{\sqrt 5-1}{4}.
      \]
  \end{itemize}
\end{tcolorbox}

\begin{significationintuitive}
Un cosinus n'est rien d'autre que l'abscisse d'un point du cercle unité.
Les angles « remarquables » sont ceux dont le point correspondant possède des
coordonnées \emph{algébriques simples} — exprimables par des racines carrées
emboîtées. Derrière ce confort calculatoire se cache la théorie des polygones
réguliers constructibles à la règle et au compas~: le carré ($\pi/4$),
l'hexagone ($\pi/3$) et le pentagone ($\pi/5$) le sont, ce qui explique
précisément pourquoi leurs cosinus s'écrivent par radicaux.
\end{significationintuitive}

\decsep{}

%=============================================================
%  § 2 — DÉMONSTRATION
%=============================================================
\secnum{navyblue}{2}{Démonstration}

\begin{ideepreuve}
Les angles $\pi/4$, $\pi/3$, $\pi/6$ se lisent sur deux triangles classiques.
À partir de là, les formules d'addition et de bissection engendrent $\pi/12$ et
$\pi/8$. Le cas $\pi/5$ est de nature différente~: il faut exploiter une relation
algébrique entre $\cos\theta$ et $\cos(n\theta)$ pour obtenir une équation
polynomiale, dont la bonne racine donne la valeur cherchée.
\end{ideepreuve}

\vspace{10pt}
\rubriq{Preuve}

\begin{mdframed}[style=proofbar]
  \textit{Étape 1 — L'angle $\pi/4$ par symétrie.}
  Dans un triangle rectangle isocèle, les deux angles aigus valent $\pi/4$,
  donc $\cos(\pi/4)=\sin(\pi/4)$. L'identité $\cos^2+\sin^2=1$ donne
  $2\cos^2(\pi/4)=1$, d'où $\cos(\pi/4)=\dfrac{\sqrt 2}{2}$ (signe positif car
  $\pi/4\in\,]0,\pi/2[$).

  \medskip
  \textit{Étape 2 — Les angles $\pi/3$ et $\pi/6$ par le triangle équilatéral.}
  Un triangle équilatéral de côté $1$ a tous ses angles égaux à $\pi/3$. La
  hauteur le coupe en deux triangles rectangles d'hypoténuse $1$ et de petit côté
  $1/2$, l'angle adjacent à ce petit côté étant $\pi/3$~: donc
  $\cos(\pi/3)=\dfrac{1/2}{1}=\dfrac{1}{2}$. Par Pythagore la hauteur vaut
  $\sqrt{1-\tfrac14}=\dfrac{\sqrt 3}{2}$, d'où $\sin(\pi/3)=\dfrac{\sqrt 3}{2}$.
  Comme $\pi/6=\pi/2-\pi/3$, les formules de complément donnent
  $\cos(\pi/6)=\sin(\pi/3)=\dfrac{\sqrt 3}{2}$ et $\sin(\pi/6)=\dfrac12$.

  \medskip
  \textit{Étape 3 — $\pi/12$ par différence d'angles.}
  Puisque $\dfrac{\pi}{12}=\dfrac{\pi}{3}-\dfrac{\pi}{4}$,
  \[
    \cos\frac{\pi}{12}
    =\cos\frac{\pi}{3}\cos\frac{\pi}{4}+\sin\frac{\pi}{3}\sin\frac{\pi}{4}
    =\frac12\cdot\frac{\sqrt2}{2}+\frac{\sqrt3}{2}\cdot\frac{\sqrt2}{2}
    =\frac{\sqrt 6+\sqrt 2}{4}.
  \]

  \medskip
  \textit{Étape 4 — $\pi/8$ par bissection.}
  De $\dfrac{\pi}{8}=\dfrac12\cdot\dfrac{\pi}{4}$ et de la formule de bissection
  (signe positif car $\pi/8\in\,]0,\pi/2[$)~:
  \[
    \cos\frac{\pi}{8}=\sqrt{\frac{1+\cos(\pi/4)}{2}}
    =\sqrt{\frac{1+\tfrac{\sqrt2}{2}}{2}}
    =\frac{\sqrt{2+\sqrt 2}}{2}.
  \]

  \medskip
  \textit{Étape 5 — $\pi/5$ par une équation polynomiale.}
  Posons $\theta=\dfrac{\pi}{5}$ et $x=\cos\theta$. Comme $5\theta=\pi$, on a
  $3\theta=\pi-2\theta$, donc $\cos(3\theta)=-\cos(2\theta)$. En développant
  $\cos(3\theta)=4x^3-3x$ et $\cos(2\theta)=2x^2-1$~:
  \[
    4x^3-3x=-(2x^2-1)
    \;\Longleftrightarrow\;
    4x^3+2x^2-3x-1=0.
  \]
  Le réel $x=-1$ (correspondant à $\theta=\pi$) est racine évidente~; la
  division euclidienne factorise
  \[
    4x^3+2x^2-3x-1=(x+1)(4x^2-2x-1).
  \]
  Le facteur $4x^2-2x-1$ a pour racines $\dfrac{1\pm\sqrt5}{4}$. Comme
  $\cos(\pi/5)>0$, on retient
  \[
    \cos\frac{\pi}{5}=\frac{1+\sqrt 5}{4}=\frac{\varphi}{2},
  \]
  où $\varphi=\dfrac{1+\sqrt5}{2}$ est le \emph{nombre d'or}. La racine négative
  $\dfrac{1-\sqrt5}{4}$ vaut $\cos(3\pi/5)=-\cos(2\pi/5)$, d'où
  $\cos\dfrac{2\pi}{5}=\dfrac{\sqrt5-1}{4}$.
  \hfill$\blacksquare$
\end{mdframed}

\decsep{}

%=============================================================
%  § 3 — EXEMPLES, CONTRE-EXEMPLES ET EXERCICES
%=============================================================
\secnum{exgreen}{3}{Exemples, contre-exemples et exercices}

\rubriq{Exemples}

\begin{mdframed}[style=exbar]
  \textbf{1. Cosinus emboîtés.}
  La bissection répétée de $\pi/4$ donne
  \[
    \cos\frac{\pi}{8}=\frac{\sqrt{2+\sqrt2}}{2},\qquad
    \cos\frac{\pi}{16}=\frac{\sqrt{2+\sqrt{2+\sqrt2}}}{2},
  \]
  motif qui se prolonge indéfiniment (lié à la formule de Viète pour $\pi$).

  \smallskip\noindent
  \textbf{2. La famille du pentagone.}
  À partir de $\cos(\pi/5)$ et d'une bissection,
  $\cos\dfrac{\pi}{10}=\cos 18^\circ=\dfrac{\sqrt{10+2\sqrt5}}{4}$.

  \smallskip\noindent
  \textbf{3. Somme et différence réunies.}
  $\cos\dfrac{5\pi}{12}=\cos\!\big(\tfrac{\pi}{4}+\tfrac{\pi}{6}\big)
   =\dfrac{\sqrt6-\sqrt2}{4}$, complément exact de $\cos(\pi/12)$.

  \smallskip\noindent
  \textbf{4. Lien avec les autres leçons.}
  Les nombres $\cos\!\big(\tfrac{2k\pi}{n}\big)$ sont les parties réelles des
  racines $n$-èmes de l'unité~: ce sont des racines des polynômes de
  Tchebychev et leur degré algébrique est régi par les polynômes cyclotomiques.
\end{mdframed}

\vspace{6pt}
\rubriq{Contre-exemples et pièges}

\begin{mdframed}[style=cexbar]
  \textbf{Tous les angles ne sont pas « remarquables ».}\enspace%
  $\cos(\pi/7)$ et $\cos(\pi/9)=\cos 20^\circ$ \emph{ne s'expriment pas} par
  radicaux réels emboîtés~: l'heptagone et l'ennéagone réguliers ne sont pas
  constructibles à la règle et au compas (angle non trisectable). De
  $\cos\!\big(3\cdot\tfrac{\pi}{9}\big)=\cos\tfrac{\pi}{3}=\tfrac12$ et de
  $\cos 3\theta=4\cos^3\theta-3\cos\theta$, on tire seulement l'équation
  cubique irréductible $8x^3-6x-1=0$ pour $\cos(\pi/9)$, sans solution
  « propre » par radicaux réels (\textit{casus irreducibilis}).

  \smallskip\noindent
  \textbf{Erreur de signe à la bissection.}\enspace%
  La formule $\cos(\theta/2)=\pm\sqrt{(1+\cos\theta)/2}$ exige de choisir le
  signe selon le \emph{quadrant} de $\theta/2$~: positif sur $]0,\pi/2[$,
  négatif sur $]\pi/2,\pi[$. L'oublier mène à $\cos(3\pi/4)=+\tfrac{\sqrt2}{2}$,
  ce qui est faux.

  \smallskip\noindent
  \textbf{$\pi/5$ n'est pas $36$ « tout court ».}\enspace%
  Travailler en degrés sans conversion, ou confondre $\cos(\pi/5)$ et
  $\cos(2\pi/5)$, fait perdre le facteur d'échelle~: $\cos(\pi/5)$ et
  $\cos(3\pi/5)=-\cos(2\pi/5)$ sont les deux racines de $4x^2-2x-1=0$,
  tandis que $\cos(2\pi/5)$ est, lui, racine de $4x^2+2x-1=0$
  (cf.\ exercice~3).
\end{mdframed}

\vspace{6pt}
\exolabel

\begin{mdframed}[style=exobar]
  \begin{enumerate}
    \item Calculer $\cos\dfrac{7\pi}{12}$ et $\sin\dfrac{\pi}{12}$ de deux
          manières (différence d'angles, puis $\cos^2+\sin^2=1$).

    \item Montrer, sans calculer chaque facteur, que
          $\cos\dfrac{\pi}{8}\,\cos\dfrac{3\pi}{8}=\dfrac{\sqrt2}{4}$.
          \textit{Indication}~: $\cos\dfrac{3\pi}{8}=\sin\dfrac{\pi}{8}$, puis
          la formule de duplication $2\sin\alpha\cos\alpha=\sin 2\alpha$.

    \item Reprendre la méthode de l'étape~5 pour montrer
          $4\cos^2\dfrac{2\pi}{5}+2\cos\dfrac{2\pi}{5}-1=0$ et retrouver
          $\cos\dfrac{2\pi}{5}=\dfrac{\sqrt5-1}{4}$.

    \item Vérifier les identités remarquables
          $\cos\dfrac{\pi}{5}-\cos\dfrac{2\pi}{5}=\dfrac12$ et
          $\cos\dfrac{\pi}{5}\,\cos\dfrac{2\pi}{5}=\dfrac14$. En déduire que
          $\cos\dfrac{\pi}{5}$ et $\cos\dfrac{2\pi}{5}$ sont les racines de
          $t^2-\dfrac{\sqrt5}{2}\,t+\dfrac14=0$, et retrouver ainsi leurs valeurs.

    \item (Produit de Viète) Montrer que
          $\dfrac{2}{\pi}=\cos\dfrac{\pi}{4}\cdot\cos\dfrac{\pi}{8}\cdot
          \cos\dfrac{\pi}{16}\cdots$ en itérant $\sin\theta=2\sin\tfrac\theta2\cos\tfrac\theta2$.

    \item Montrer que $\cos\dfrac{\pi}{7}$ est racine de $8x^3-4x^2-4x+1=0$~;
          en déduire $\cos\dfrac{\pi}{7}+\cos\dfrac{3\pi}{7}+\cos\dfrac{5\pi}{7}=\dfrac12$
          \textit{(somme des racines)}.
  \end{enumerate}
\end{mdframed}

\leconfooter{Trigonométrie classique ; C.\ F.\ Gauss, polygones constructibles (1796)}
